À l'issue de la deuxième année, d'ici septembre, notre approche de rejeu de
spécifications paramétriques à partir de scripts de règles doit pouvoir être
présentée à la communauté française d'informatique graphique. %
Faute de temps, nous n'aborderons pas l'animation d'un modèles sur la dimension
temps et poursuivrons alors avec l'annotation des modèles avec des données
historiques. %
Dans cette partie là, il s'agira d'associer un certains nombre de données
d'origines et de formes variées aux modèles 3D. %
Ceux-ci pourront alors, par exemple, détenir des information historiques comme
des témoignages historiques ou encore des données plutôt architecturales telles
que la textures et/ou la rugosité des murs extérieurs d'un bâtiment. %
Cette partie là sera, dès mars 2024, directement suivie de la rédaction du
manuscript de thèse en vue de soutenir à la fin de cette même année. %


% D'ici décembre, nous prévoyons d'intégrer les cas de fusion, d'implémenter
% l'édition des spécifications paramétriques avec l'ajout, la suppression et le
% déplacement d'appels de règles. %
% Puis nous enchaînerons en fin d'année sur l'implémentation du rejeu de
% spécifications paramétriques après édition. %
% Les arbres d'appariement permettront ainsi de tracer les évènements et de
% proposer des options dans les cas où plusieurs paramètres topologiques
% apparaissent ou si, au contraire, aucun paramètre topologique ne peut être
% trouvé. %
% L'intégration d'une dimension temporelle pour animer l'évolution d'un
% modèle devra certainement être mise de côté par soucis de temps afin d'une
% part d'étendre notre travail à l'utilisation d'opérations plus complexes
% définies à partir de scripts de règles, et d'autre part d'intégrer au modèle des
% annotations historiques. %


%%% Local Variables:
%%% mode: latex
%%% TeX-master: "../main"
%%% End:
